\section{Data Preparation}

The data preparation phase covers all activities to construct the final dataset for modeling from the initial raw data. The implementation materializes two parallel modeling exports (30-day and 90-day label definitions) with the same feature side; the final choice of window for training belongs to the modeling phase, not to this step. Below, each step states \emph{what} is done, \emph{why} it is done, and \emph{what} remains in the data after the step.

\subsection{Select Data}

\subsubsection{Data selection during extraction}

\paragraph{What was already fixed in the warehouse:}
\begin{itemize}
  \item \textbf{Schema focus:} the GitHub Archive extract keeps identifiers (repository, actor), event type and time, repository URL, visibility, and a derived \texttt{season} label; large or rarely needed structures (e.g.\ deep payload and most nested metadata) are excluded to limit size and to align rows across tables.
  \item \textbf{Event types:} the pipeline uses \texttt{ForkEvent} and \texttt{PullRequestEvent} for the conversion target, and \texttt{WatchEvent}, \texttt{PushEvent}, \texttt{IssuesEvent}, \texttt{IssueCommentEvent}, and \texttt{ReleaseEvent} for activity features, with the seasonal windows and repository filters set in the collection phase.
  \item \textbf{Why:} a narrow, consistent schema reduces noise, keeps joins predictable, and matches the business question (activity and fork-to-PR behaviour within defined horizons).
\end{itemize}

\subsubsection{Data selection during EDA and modeling}

\paragraph{Columns dropped from the modeling copy of the activity table.}
\begin{itemize}
  \item \textbf{\texttt{public}:} no usable variation in this sample, so it does not add signal and is removed to avoid a constant feature.
  \item \textbf{\texttt{repo\_url}:} not used as a predictive input in the current design; the identifier \texttt{repo\_name} suffices. The URL can be reattached later if static metadata is joined from outside.
\end{itemize}

\paragraph{What is retained.}
\begin{itemize}
  \item \textbf{Keys:} \texttt{repo\_name} and \texttt{season} for every repository--season row from the sample.
  \item \textbf{Activity counters:} watches, pushes, issues, issue comments, releases (five nonnegative count variables).
  \item \textbf{Copy used in code:} \texttt{prep\_activity} is the activity table after the drops above, so all subsequent cleaning and feature steps apply to this object only.
\end{itemize}

\paragraph{Row retention.}
\begin{itemize}
  \item No repository--season rows are removed at this step: the keys from the extract are kept.
  \item Alignment with labels is deferred to a later inner join, so any inconsistency is resolved explicitly rather than by silent row drops here.
\end{itemize}

\subsection{Clean Data}

\paragraph{Principle.}
\begin{itemize}
  \item Activity metrics show heavy tails (very popular repositories). These are treated as legitimate extremes, not as automatic row-level errors, so the strategy is tail reduction, not deletion of outlying repositories.
  \item Fork and PR event lists used for labels are not winsorized: labeling relies on raw timestamps and user linkage; only tabular activity counts are compressed for modeling stability.
\end{itemize}

\paragraph{Winsorization of activity counts.}
For each of the five activity counters:
\begin{enumerate}
  \item Values are cast to \texttt{float64} before quantiles, so nullable integer types from Parquet do not break compatibility with real-valued quantile thresholds.
  \item The $0.99$ quantile is computed (fixed as \texttt{WINSOR\_Q = 0.99} in the implementation).
  \item Values above that quantile are capped; capped values are written to \texttt{*\_winsor} columns; original columns are retained unchanged.
\end{enumerate}

\paragraph{Outcome of cleaning.}
The modeling pipeline can use either raw or capped scales where appropriate; the full sample and event-level label inputs remain available without arbitrary row loss.


\subsection{Construct Data}

\subsubsection{Sampling strategy}

\begin{itemize}
  \item \textbf{Unit of analysis:} one row = one repository in one season (not a single row aggregated across the whole year for the same repo).
  \item \textbf{Sample size:} defined upstream (e.g.\ cap and filters in the extract); this phase does not re-subsample for class balance. The tables exported after merge are the full set of keys that pass join and label construction.
  \item \textbf{Why:} preserves the design of the initial draw and makes train/validation splits in modeling transparent relative to that population.
\end{itemize}

\subsubsection{Matching forks to pull requests (label definition)}

\paragraph{Conversion rule (per fork).}
\begin{itemize}
  \item A fork \textbf{converts} if there is at least one \texttt{PullRequestEvent} such that: (i) \texttt{user\_login} matches the fork author, (ii) PR time is strictly after the fork time, (iii) PR time falls within a chosen observation window (in calendar days) after the fork.
  \item PRs on or before the fork, or after the end of the window, do not count toward conversion for that fork.
\end{itemize}

\paragraph{Two observation windows (not a single default).}
The implementation sets\\ \texttt{OBSERVATION\_WINDOWS\_DAYS = (30, 90)} and runs the same labeling logic for each value. There is no single default observation length inside the data-preparation code: the notebook builds \texttt{labels\_30d} and \texttt{labels\_90d} and writes
\texttt{fork\_conversion\_labels\_obs30d} and
\texttt{fork\_conversion\_labels\_obs90d} (Parquet, or CSV if no Parquet engine is available). Each row includes \texttt{observation\_window\_days} (30 or 90).

\paragraph{Sensitivity of ranks between windows.}
A separate check compares repository-level\\ \texttt{fork\_conversion\_rate} at 30 and 90~days (Spearman/Kendall correlation, decile reversals, and the share of 90-day conversion events not attributable within 30~days) and prints textual guidance for modeling (which file may be a reasonable first choice). That guidance is for the modeling step only; it does not change the fact that both full modeling tables are produced in Format Data.

\paragraph{Aggregates per repository and season.}
For each \texttt{repo\_name} and \texttt{season}, the construction step records:
\begin{itemize}
  \item \texttt{n\_forks} --- number of fork events in the cell;
  \item \texttt{n\_converted\_forks} --- forks that meet the conversion rule within the chosen window;
  \item \texttt{fork\_conversion\_rate} = converted forks / fork count (this is not the same as total PR events divided by forks without user--time matching);
  \item \texttt{n\_pr\_events} --- count of PR events in the list (for diagnostics and scale).
\end{itemize}

\subsubsection{Derived features (activity side)}

\begin{itemize}
  \item For each of the five activity variables, log1p is applied to: (a) the raw count, and (b) the winsorized (\texttt{*\_winsor}) value, producing \texttt{log1p\_*} features suited to strong skew and nonnegative support.
  \item Original levels and \texttt{*\_winsor} columns remain in the table so the analyst can choose scale or run robustness checks later.
\end{itemize}

\subsubsection{Generated records}

\begin{itemize}
  \item \textbf{Grain:} one flat row per \texttt{repo\_name} and \texttt{season}; nested fork/PR arrays are not written to these flat label files.
  \item \textbf{Empty PR lists:} rows are kept; conversion can be zero; such observations are not removed.
  \item \textbf{Only} fork--PR pairs that satisfy the match rules and fall inside the active window update \texttt{n\_converted\_forks} and the rate.
\end{itemize}

\subsection{Integrate Data}

\paragraph{Join logic.}
\begin{itemize}
  \item The prepared activity frame (\texttt{prep\_activity}: identifiers, counts, winsorized and log-transformed fields) is merged with a label table on \texttt{repo\_name} and \texttt{season}.
  \item The join is an inner merge with a one-to-one key check (\texttt{validate="one\_to\_one"}), and row counts are asserted equal to \texttt{prep\_activity} and the label table. Each key appears at most once on each side and no duplicate fan-out is introduced.
  \item \textbf{Two integrated frames, not one ``default'':} the code builds \texttt{modeling\_df\_30d} and\\ \texttt{modeling\_df\_90d} by merging to \texttt{labels\_30d} and \texttt{labels\_90d} respectively. Both are present for the same set of keys in the prepared sample; the notebook does not pick a single ``main'' target here---that choice is deferred to the modeling stage when loading
  \texttt{modeling\_dataset\_obs30d} or
  \texttt{modeling\_dataset\_obs90d}.
  \item \textbf{Result (per window):} one wide row per repository and season, combining feature columns with label fields (\texttt{observation\_window\_days}, \texttt{n\_forks}, \texttt{n\_converted\_forks},\\ \texttt{fork\_conversion\_rate}, \texttt{n\_pr\_events}, and related fields from construction).
\end{itemize}

\subsection{Format Data}

\paragraph{Purpose.} Stabilize column order, add convenient identifiers and optional stratification helpers, and export artifacts for the modeling phase.

\paragraph{Per-window merge and export pipeline.}
The export step applies, for each label table, an inner join of \texttt{prep\_activity} to the chosen labels (same key logic as in Integrate Data), then:
\begin{itemize}
  \item \texttt{repo\_id}: integer code for \texttt{repo\_name} via categorical codes; \texttt{repo\_name} is retained.
  \item \texttt{has\_conversion}: $1$ if \texttt{n\_converted\_forks} $>0$, else $0$.
  \item \texttt{activity\_strength}: sum of the five \texttt{log1p\_} features on raw activity counts.
  \item \texttt{stratum\_activity\_quintile}, \texttt{stratum\_fork\_quintile}: quintile bins (0--4) from \texttt{pd.qcut} on \texttt{activity\_strength} and on \texttt{n\_forks} (\texttt{duplicates="drop"} on degenerate inputs), for optional stratified splits---they need not be used as model features.
  \item \textbf{Column order} in the exported file: meta (\texttt{repo\_name}, \texttt{repo\_id}, \texttt{season}), then all \texttt{log1p\_*} columns (sorted by name in code), then the stratum block above, then the label block\\ \texttt{observation\_window\_days}, \texttt{n\_forks}, \texttt{n\_converted\_forks}, \texttt{fork\_conversion\_rate},\\ \texttt{n\_pr\_events}, \texttt{has\_conversion}, then remaining columns (raw counts, \texttt{*\_winsor}, etc.).
  \item \textbf{Shuffle:} the merged table is permuted with \texttt{sample(frac=1.0, random\_state=335)} so the file order is not the extract order, while runs remain reproducible.
\end{itemize}

\paragraph{Output files.}
\begin{itemize}
  \item \textbf{Label-only files (one per window):} \texttt{fork\_conversion\_labels\_obs30d} and\\ \texttt{fork\_conversion\_labels\_obs90d} in Parquet or CSV, as written in the construction step.
  \item \textbf{Full modeling tables (one per window):} \texttt{modeling\_dataset\_obs30d} and\\ \texttt{modeling\_dataset\_obs90d} (Parquet, or \texttt{.csv} with the same stem if no Parquet engine is available). They share the same feature and metadata layout; only the label columns differ. 
\end{itemize}
